\chapter{Of the Society which was Formed in the Rue des Lombards}

\lettrine{A}{fter} a moment's silence, in which D'Artagnan appeared to be collecting, not one idea but all his ideas,—<It cannot be, my dear Planchet,> said he, <that you have not heard of his majesty \regnal{Charles I} of England?>

\zz
<Alas! yes, monsieur, since you left France in order to assist him, and that, in spite of that assistance, he fell, and was near dragging you down in his fall.>

<Exactly so; I see you have a good memory, Planchet.>

<\swear{Peste!} the astonishing thing would be, if I could have lost that memory, however bad it might have been. When one has heard Grimaud, who, you know, is not given to talking, relate how the head of King Charles fell, how you sailed the half of a night in a scuttled vessel, and saw floating on the water that good M.~Mordaunt with a certain gold-hafted dagger buried in his breast, one is not very likely to forget such things.>

<And yet there are people who forget them, Planchet.>

<Yes, such as have not seen them, or have not heard Grimaud relate them.>

<Well, it is all the better that you recollect all that; I shall only have to remind you of one thing, and that is that \regnal{Charles I} had a son.>

<Without contradicting you, monsieur, he had two,> said Planchet; <for I saw the second one in Paris, M.~le Duke of York, one day, as he was going to the Palais Royal, and I was told that he was not the eldest son of \regnal{Charles I}. As to the eldest, I have the honour of knowing him by name, but not personally.>

<That is exactly the point, Planchet, we must come to: it is to this eldest son, formerly called the Prince of Wales, and who is now styled \regnal{Charles II}, king of England.>

<A king without a kingdom, monsieur,> replied Planchet, sententiously.

<Yes, Planchet, and you may add an unfortunate prince, more unfortunate than the poorest man of the people lost in the worst quarter of Paris.>

Planchet made a gesture full of that sort of compassion which we grant to strangers with whom we think we can never possibly find ourselves in contact. Besides, he did not see in this politico-sentimental operation any sign of the commercial idea of M.~d'Artagnan, and it was in this idea that D'Artagnan, who was, from habit, pretty well acquainted with men and things, had principally interested Planchet.

<I am come to our business. This young Prince of Wales, a king without a kingdom, as you have so well said, Planchet, has interested me. I, D'Artagnan, have seen him begging assistance of Mazarin, who is a miser, and the aid of Louis, who is a child, and it appeared to me, who am acquainted with such things, that in the intelligent eye of the fallen king, in the nobility of his whole person, a nobility apparent above all his miseries, I could discern the stuff of a man and the heart of a king.>

Planchet tacitly approved of all this; but it did not at all, in his eyes at least, throw any light upon D'Artagnan's idea. The latter continued: <This, then, is the reasoning which I made with myself. Listen attentively, Planchet, for we are coming to the conclusion.>

<I am listening.>

<Kings are not so thickly sown upon the earth, that people can find them whenever they want them. Now, this king without a kingdom is, in my opinion, a grain of seed which will blossom in some season or other, provided a skilful, discreet, and vigorous hand sow it duly and truly, selecting soil, sky, and time.>

Planchet still approved by a nod of his head, which showed that he did not perfectly comprehend all that was said.

<<Poor little seed of a king,> said I to myself, and really I was affected, Planchet, which leads me to think I am entering upon a foolish business. And that is why I wished to consult you, my friend.>

Planchet coloured with pleasure and pride.

<<Poor little seed of a king! I will pick you up and cast you into good ground.>>

<Good God!> said Planchet, looking earnestly at his old master, as if in doubt as to the state of his reason.

<Well, what is it?> said D'Artagnan; <who hurts you?>

<Me! nothing, monsieur.>

<You said, <Good God!>>

<Did I\@?>

<I am sure you did. Can you already understand?>

<I confess, M.~d'Artagnan, that I am afraid\longdash>

<To understand?>

<Yes.>

<To understand that I wish to replace upon his throne this King \regnal{Charles II}, who has no throne? Is that it?>

Planchet made a prodigious bound in his chair. <Ah, ah!> said he, in evident terror, <that is what you call a restoration!>

<Yes, Planchet; is it not the proper term for it?>

<Oh, no doubt, no doubt! But have you reflected seriously?>

<Upon what?>

<Upon what is going on yonder.>

<Where?>

<In England.>

<And what is that? Let us see, Planchet.>

<In the first place, monsieur, I ask you pardon for meddling in these things, which have nothing to do with my trade; but since it is an affair that you propose to me—for you are proposing an affair, are you not?—>

<A superb one, Planchet.>

<But as it is business you propose to me, I have the right to discuss it.>

<Discuss it, Planchet; out of discussion is born light.>

<Well, then, since I have monsieur's permission, I will tell him that there is yonder, in the first place, the parliament.>

<Well, next?>

<And then the army.>

<Good! Do you see anything else?>

<Why, then the nation.>

<Is that all?>

<The nation which consented to the overthrow and death of the late king, the father of this one, and which will not be willing to belie its acts.>

<Planchet,> said D'Artagnan, <you argue like a cheese! The nation—the nation is tired of these gentlemen who give themselves such barbarous names, and who sing songs to it. Chanting for chanting, my dear Planchet; I have remarked that nations prefer singing a merry chant to the plain chant. Remember the Fronde; what did they sing in those times? Well, those were good times.>

<Not too good, not too good! I was near being hung in those times.>

<Well, but you were not.>

<No.>

<And you laid the foundations of your fortune in the midst of all those songs?>

<That is true.>

<Then you have nothing to say against them.>

<Well, I return, then, to the army and parliament.>

<I say that I borrow twenty thousand livres of M.~Planchet, and that I put twenty thousand livres of my own to it; and with these forty thousand livres I raise an army.>

Planchet clasped his hands; he saw that D'Artagnan was in earnest, and, in good truth, he believed his master had lost his senses.

<An army!—ah, monsieur,> said he, with his most agreeable smile, for fear of irritating the madman, and rendering him furious,—<an army!—how many?>

<Of forty men,> said D'Artagnan.

<Forty against forty thousand! that is not enough. I know very well that you, M.~d'Artagnan, alone, are equal to a thousand men; but where are we to find thirty-nine men equal to you? Or, if we could find them, who would furnish you with money to pay them?>

<Not bad, Planchet. Ah, the devil! you play the courtier.>

<No, monsieur, I speak what I think, and that is exactly why I say that, in the first pitched battle you fight with your forty men, I am very much afraid\longdash>

<Therefore I shall fight no pitched battles, my dear Planchet,> said the Gascon, laughing. <We have very fine examples in antiquity of skilful retreats and marches, which consisted in avoiding the enemy instead of attacking them. You should know that, Planchet, you who commanded the Parisians the day on which they ought to have fought against the musketeers, and who so well calculated marches and countermarches, that you never left the Palais Royal.>

Planchet could not help laughing. <It is plain,> replied he, <that if your forty men conceal themselves, and are not unskilful, they may hope not to be beaten: but you propose obtaining some result, do you not?>

<No doubt. This, then, in my opinion, is the plan to be proceeded upon in order quickly to replace his majesty \regnal{Charles II} on his throne.>

<Good!> said Planchet, increasing his attention; <let us see your plan. But in the first place it seems to me we are forgetting something.>

<What is that?>

<We have set aside the nation, which prefers singing merry songs to psalms, and the army, which we will not fight; but the parliament remains, and that seldom sings.>

<Nor does it fight. How is it, Planchet, that an intelligent man like yourself should take any heed of a set of brawlers who call themselves Rumps and Barebones? The parliament does not trouble me at all, Planchet.>

<As soon as it ceases to trouble you, monsieur, let us pass on.>

<Yes, and arrive at the result. You remember Cromwell, Planchet?>

<I have heard a great deal of talk about him.>

<He was a rough soldier.>

<And a terrible eater, moreover.>

<What do you mean by that?>

<Why, at one gulp he swallowed all England.>

<Well, Planchet, the evening before the day on which he swallowed England, if any one had swallowed M.~Cromwell?>

<Oh, monsieur, it is one of the axioms of mathematics that the container must be greater than the contained.>

<Very well! That is our affair, Planchet.>

<But M.~Cromwell is dead, and his container is now the tomb.>

<My dear Planchet, I see with pleasure that you have not only become a mathematician, but a philosopher.>

<Monsieur, in my grocery business I use much printed paper, and that instructs me.>

<Bravo! You know then, in that case—for you have not learnt mathematics and philosophy without a little history—that after this Cromwell so great, there came one who was very little.>

<Yes; he was named Richard, and he as done as you have, M.~ d'Artagnan—he has tendered his resignation.>

<Very well said—very well! After the great man who is dead, after the little one who tendered his resignation, there came a third. This one is named Monk; he is an able general, considering he has never fought a battle; he is a skilful diplomatist, considering that he never speaks in public, and that having to say <good-day> to a man, he meditates twelve hours, and ends by saying <good night;> which makes people exclaim <miracle!> seeing that it falls out correctly.>

<That is rather strong,> said Planchet; <but I know another political man who resembles him very much.>

<M.~Mazarin you mean?>

<Himself.>

<You are right, Planchet; only M.~Mazarin does not aspire to the throne of France; and that changes everything. Do you see? Well, this M.~Monk, who has England ready-roasted in his plate, and who is already opening his mouth to swallow it—this M.~Monk, who says to the people of \regnal{Charles II}, and to \regnal{Charles II} himself, <Nescio vos>—>

<I don't understand English,> said Planchet.

<Yes, but I understand it,> said D'Artagnan. <<Nescio vos> means <I do not know you.> This M.~ Monk, the most important man in England, when he shall have swallowed it\longdash>

<Well?> asked Planchet.

<Well, my friend, I shall go over yonder, and with my forty men I shall carry him off, pack him up, and bring him into France, where two modes of proceeding present themselves to my dazzled eyes.>

<Oh! and to mine too,> cried Planchet, transported with enthusiasm. <We will put him in a cage and show him for money.>

<Well, Planchet, that is a third plan, of which I had not thought.>

<Do you think it a good one?>

<Yes, certainly, but I think mine better.>

<Let us see yours, then.>

<In the first place, I shall set a ransom on him.>

<Of how much?>

<\swear{Peste!} a fellow like that must be well worth a hundred thousand crowns.>

<Yes, yes!>

<You see, then—in the first place, a ransom of a hundred thousand crowns.>

<Or else\longdash>

<Or else, what is much better, I deliver him up to King Charles, who, having no longer either a general or an army to fear, nor a diplomatist to trick him, will restore himself, and when once restored, will pay down to me the hundred thousand crowns in question. That is the idea I have formed; what do you say to it, Planchet?>

<Magnificent, monsieur!> cried Planchet, trembling with emotion. <How did you conceive that idea?>

<It came to me one morning on the banks of the Loire, whilst our beloved king, \regnal{Louis XIV}, was pretending to weep upon the hand of Mademoiselle de Mancini.>

<Monsieur, I declare the idea is sublime. But\longdash>

<Ah! is there a but?>

<Permit me! But this is a little like the skin of that fine bear—you know—that they were about to sell, but which it was necessary to take from the back of the living bear. Now, to take M.~Monk, there will be a bit of a scuffle, I should think.>

<No doubt; but as I shall raise an army to\longdash>

<Yes, yes—I understand, \swear{parbleu!}—a coup-de-main. Yes, then, monsieur, you will triumph, for no one equals you in such sorts of encounters.>

<I certainly am lucky in them,> said D'Artagnan, with a proud simplicity. <You know that if for this affair I had my dear Athos, my brave Porthos, and my cunning Aramis, the business would be settled; but they are all lost, as it appears, and nobody knows where to find them. I will do it, then, alone. Now, do you find the business good, and the investment advantageous?>

<Too much so—too much so.>

<How can that be?>

<Because fine things never reach the expected point.>

<This is infallible, Planchet, and the proof is that I undertake it. It will be for you a tolerably pretty gain, and for me a very interesting stroke. It will be said, <Such was the old age of M.~d'Artagnan,> and I shall hold a place in tales and even in history itself, Planchet. I am greedy of honour.>

<Monsieur,> cried Planchet, <when I think that it is here, in my home, in the midst of my sugar, my prunes, and my cinnamon, that this gigantic project is ripened, my shop seems a palace to me.>

<Beware, beware, Planchet! If the least report of this escapes, there is the Bastille for both of us. Beware, my friend, for this is a plot we are hatching. M.~Monk is the ally of M.~Mazarin—beware!>

<Monsieur, when a man has had the honour to belong to you, he knows nothing of fear; and when he has had the advantage of being bound up in interests with you, he holds his tongue.>

<Very well; that is more your affair than mine, seeing that in a week I shall be in England.>

<Depart, monsieur, depart—the sooner the better.>

<Is the money, then, ready?>

<It will be to-morrow; to-morrow you shall receive it from my own hands. Will you have gold or silver?>

<Gold; that is most convenient. But how are we going to arrange this? Let us see.>

<Oh, good Lord! in the simplest way possible. You shall give me a receipt, that is all.>

<No, no,> said D'Artagnan, warmly; <we must preserve order in all things.>

<That is likewise my opinion; but with you, M.~ d'Artagnan\longdash>

<And if I should die yonder—if I should be killed by a musket-ball—if I should burst from drinking beer?>

<Monsieur, I beg you to believe that in that case I should be so much afflicted at your death, that I should not think about the money.>

<Thank you, Planchet; but no matter. We shall, like two lawyers' clerks, draw up together an agreement, a sort of act, which may be called a deed of company.>

<Willingly, monsieur.>

<I know it is difficult to draw such a thing up, but we can try.>

<Let us try, then.> And Planchet went in search of pens, ink, and paper. D'Artagnan took the pen and wrote:—<Between Messire d'Artagnan, ex-lieutenant of the king's musketeers, at present residing in the Rue Tiquetonne, Hôtel de la Chevrette; and the Sieur Planchet, grocer, residing in the Rue des Lombards, at the sign of the \buildingname{Pilon d'Or,} it has been agreed as follows:—A company, with a capital of forty thousand livres, and formed for the purpose of carrying out an idea conceived by M.~d'Artagnan, and the said Planchet approving of it in all points, will place twenty thousand livres in the hands of M.~d'Artagnan. He will require neither repayment nor interest before the return of M.~d'Artagnan from a journey he is about to take into England. On his part, M.~ d'Artagnan undertakes it to find twenty thousand livres, which he will join to the twenty thousand already laid down by the Sieur Planchet. He will employ the said sum of forty thousand livres according to his judgment in an undertaking which is described below. On the day when M.~d'Artagnan shall have re-established, by whatever means, his majesty King \regnal{Charles II} upon the throne of England, he will pay into the hands of M.~Planchet the sum of\longdash>

<The sum of a hundred and fifty thousand livres,> said Planchet, innocently, perceiving that D'Artagnan hesitated.

<Oh, the devil, no!> said D'Artagnan, <the division cannot be made by half; that would not be just.>

<And yet, monsieur, we each lay down half,> objected Planchet, timidly.

<Yes; but listen to this clause, my dear Planchet, and if you do not find if equitable in every respect when it is written, well, we can scratch it out again:—<Nevertheless, as M.~d'Artagnan brings to the association, besides his capital of twenty thousand livres, his time, his idea, his industry, and his skin,—things which he appreciates strongly, particularly the last,—M.~d'Artagnan will keep, of the three hundred thousand livres, two hundred thousand livres for himself, which will make his share two-thirds.>>

<Very well,> said Planchet.

<Is it just?> asked D'Artagnan.

<Perfectly just, monsieur.>

<And you will be contented with a hundred thousand livres?>

<\swear{Peste!} I think so. A hundred thousand for twenty thousand!>

<And in a month, understand.>

<How, in a month?>

<Yes, I only ask one month.>

<Monsieur,> said Planchet, generously, <I give you six weeks.>

<Thank you,> replied the musketeer, politely; after which the two partners reperused their deed.

<That is perfect, monsieur,> said Planchet; <and the late M.~ Coquenard, the first husband of Madame la Baronne du Vallon, could not have done it better.>

<Do you find it so? Let us sign it then.> And both affixed their signatures.

<In this fashion,> said D'Artagnan, <I shall be under obligations to no one.>

<But I shall be under obligations to you,> said Planchet.

<No; for whatever store I set by it, Planchet, I may lose my skin yonder, and you will lose all. À propos—\swear{peste!}—that makes me think of the principal, an indispensable clause. I shall write it:—<In case of M.~d'Artagnan dying in this enterprise, liquidation will be considered made, and the Sieur Planchet will give quittance from that moment to the shade of Messire d'Artagnan for the twenty thousand livres paid by him into the hands of the said company.>>

This last clause made Planchet knit his brows a little, but when he saw the brilliant eye, the muscular hand, the supple and strong back of his associate, he regained his courage, and, without regret, he at once added another stroke to his signature. D'Artagnan did the same. Thus was drawn the first known company contract; perhaps such things have been abused a little since, both in form and principle.

<Now,> said Planchet, pouring out the last glass of Anjou wine for D'Artagnan,—<now go to sleep, my dear master.>

<No,> replied D'Artagnan; <for the most difficult part now remains to be done, and I will think over that difficult part.>

<Bah!> said Planchet; <I have such great confidence in you, M.~d'Artagnan, that I would not give my hundred thousand livres for ninety thousand livres down.>

<And devil take me if I don't think you are right!> Upon which D'Artagnan took a candle and went up to his bedroom. 